% เตรียมตอบคำถามหลังนำเสนอระบบหลังบ้าน Colorado แยกตามคนตอบ
% สร้าง PDF:  ./qna_prep/build.sh
% ตัวเลขทั้งหมดตรงกับรายงาน (report-doc/content และ 04-chapter2-diagrams-database)
\documentclass[12pt,a4paper]{article}
\usepackage[a4paper,margin=1.8cm,top=1.6cm,bottom=1.8cm]{geometry}
\usepackage{fontspec}
\usepackage{xcolor}
\usepackage[most]{tcolorbox}
\usepackage{enumitem}
\usepackage{tabularx,booktabs,array}
\usepackage{fancyhdr}
\usepackage{pifont}

\setmainfont{Sukhumvit Set}[UprightFont={* Text}, BoldFont={* Semi Bold}]
\newfontfamily\heavy{Sukhumvit Set}[UprightFont={* Bold}]
\XeTeXlinebreaklocale "th"
\XeTeXlinebreakskip = 0pt plus 1pt
\linespread{1.25}
\setlength{\parindent}{0pt}
\setlength{\parskip}{3pt}

\definecolor{ink}{HTML}{292A2E}    % color.text
\definecolor{muted}{HTML}{505258}  % color.text.subtle
\definecolor{sA}{HTML}{FF5A36}
\definecolor{sB}{HTML}{1868DB}
\definecolor{sC}{HTML}{6E5DC6}

\newcommand{\speakerA}{นายธนภัทร พูลโภคา}
\newcommand{\speakerB}{นายอชิระ ภู่รุ่งเรือง}
\newcommand{\speakerC}{นางสาวปิยากร ผลพานิช}

\setlist[itemize]{leftmargin=1.3em, itemsep=2pt, topsep=2pt, label=\textcolor{cur}{\textbullet}}
\pagestyle{fancy}\fancyhf{}
\renewcommand{\headrulewidth}{0pt}
\fancyfoot[L]{\small\color{muted} Colorado Back-Office System  เตรียมตอบคำถาม}
\fancyfoot[R]{\small\color{muted}\thepage}

\newcommand{\speaker}[3]{% #1 สี  #2 ชื่อ  #3 หมวด
  \clearpage\colorlet{cur}{#1}%
  \begin{tcolorbox}[enhanced,colback=#1,colframe=#1,arc=4pt,left=10pt,right=10pt,top=8pt,bottom=8pt]
    \color{white}{\heavy\Huge #2}\\[2pt]{\large #3}
  \end{tcolorbox}\vspace{2pt}}
\newcounter{qa}
\newcommand{\must}{\quad{\small\textcolor{white}{\ding{72} ทุกคนต้องตอบได้}}}
\newenvironment{q}[1]{\refstepcounter{qa}%
  \begin{tcolorbox}[enhanced,colback=white,colframe=cur,boxrule=1pt,arc=4pt,breakable,
    left=8pt,right=8pt,top=4pt,bottom=6pt,before skip=6pt,after skip=6pt,
    title={\heavy Q\theqa\quad #1},colbacktitle=cur,coltitle=white,fonttitle=\large]}
  {\end{tcolorbox}}
\newcommand{\ans}{\par{\heavy\textcolor{cur}{ตอบ}}\par}
\newcommand{\evidence}[1]{\par\vspace{2pt}{\heavy\textcolor{cur}{เปิดหลักฐาน}}\quad #1\par}
\newcommand{\num}[1]{{\heavy\textcolor{cur}{#1}}}

\begin{document}
\colorlet{cur}{ink}
\begin{tcolorbox}[enhanced,colback=ink,colframe=ink,arc=4pt,left=12pt,right=12pt,top=12pt,bottom=12pt]
  \color{white}{\heavy\huge เตรียมตอบคำถาม Colorado Back-Office System}\\[4pt]
  {\large 01418321 System Analysis and Design}
\end{tcolorbox}

\vspace{6pt}
{\heavy\Large ใครตอบเรื่องอะไร}\par
\renewcommand{\arraystretch}{1.5}
\begin{tabularx}{\linewidth}{>{\bfseries}p{4.3cm} X}
\toprule
คนตอบ & หมวด\\
\midrule
\textcolor{sA}{\speakerA} & ธุรกิจ ปัญหา ตัวเลขประมาณการ ขอบเขต\\
\textcolor{sB}{\speakerB} & Business Process, Use Case, CRUD, diagram, ER\\
\textcolor{sC}{\speakerC} & SQL, ฐานข้อมูล, เว็บ, การทดสอบ, ขั้นต่อไป\\
\bottomrule
\end{tabularx}

\vspace{8pt}
{\heavy\Large ตอบ 4 จังหวะ จบใน 30 ถึง 45 วินาที}\par
\begin{enumerate}
  \item ฟังให้จบ แล้วทวนคำถามสั้น ๆ
  \item ประโยคแรกคือคำตอบ
  \item ให้หลักฐาน 1 อย่าง เช่น เปิดหน้าในรายงานหรือเปิดเว็บ
  \item โยงกลับไปที่ human-in-the-loop หรือปัญหาของเจ้าของ
\end{enumerate}

{\heavy\Large ห้ามพูด}\par
\begin{tcolorbox}[colback=ink!6,colframe=ink!6,arc=4pt]
\begin{itemize}
  \item ``เจ้าของใช้เวลา 2 ชั่วโมงต่อวัน'' แบบไม่มีคำว่าประมาณ ตัวเลขเวลาทั้งหมดเป็น\textbf{ค่าประมาณ} จากขั้นตอนที่เจ้าของเล่า
  \item ``ระบบสั่งซื้อให้อัตโนมัติ'' ที่ถูกคือระบบ\textbf{เตรียม}คำสั่งซื้อ แต่เจ้าของต้องกดอนุมัติ (5A) ก่อน
\end{itemize}
\end{tcolorbox}

% ================================================================ คนที่ 1
\speaker{sA}{\speakerA}{ธุรกิจ ปัญหา ขอบเขต}

\begin{q}{ทำไมเฟส 1 ทำแค่ Rakuten \must}
\ans
Rakuten เป็นช่องทางขายหลักที่มีคำสั่งซื้อจริงให้ทดสอบทันที และร้านใช้ RSL ซึ่งอยู่ในระบบเดียวกับ Rakuten อยู่แล้ว Rakuten RMS มี API ครอบคลุมทั้งดึงคำสั่งซื้อ ใบปะ และสต๊อก การทำทีละแพลตฟอร์มยังลดความเสี่ยงจากการเชื่อมหลายระบบพร้อมกัน
\evidence{รายงานหัวข้อ 1.1 และ 1.5}
\end{q}

\begin{q}{ตัวเลข 1 ถึง 2 ชั่วโมงต่อวันมาจากไหน \must}
\ans
มาจากจำนวนขายประมาณ \num{5 ถึง 10 ชิ้นต่อวัน} ตาม requirement คูณเวลาต่อคำสั่งซื้อประมาณ \num{12 นาที} ที่แยกเป็น 5 งาน เป็นค่าประมาณจากขั้นตอนที่เจ้าของเล่า ยังไม่ได้จับเวลาจริง ถ้าเจ้าของจับเวลาจริงให้ภายหลังจะแก้ตัวเลขในตารางที่ 1-1
\evidence{ตารางที่ 1-1 ในหัวข้อ 1.2}
\end{q}

\begin{q}{ปัญหาไม่ใช่แค่ "ไม่มีคอมพิวเตอร์" หรือ}
\ans
ไม่ใช่ เจ้าของใช้คอมพิวเตอร์และ Rakuten RMS อยู่แล้ว ปัญหาคือข้อมูลอยู่คนละที่ ต้องคัดลอกและเทียบเองทุกคำสั่งซื้อ เช่น สต๊อกอยู่ทั้งคลังบริษัทและ RSL ต้นทุนต้องรวม 8 รายการ ระบบใหม่รวมข้อมูลไว้ที่เดียวและทำขั้นตอนซ้ำให้ ปัญหาวัดผลได้จากเวลาต่อวันที่ลดลง
\end{q}

\begin{q}{ทำไมไม่ให้ระบบสั่งซื้อจาก Supplier เองเลย \must}
\ans
เพราะเป็นขั้นที่ใช้เงินจริง ถ้าข้อมูลสต๊อกหรือต้นทุนผิด ระบบอาจสั่งของผิดจำนวนโดยไม่มีใครเห็น อาจารย์แนะนำให้ใช้ human-in-the-loop เราจึงให้ระบบคำนวณต้นทุน กำไรที่คาด และเตรียมคำสั่งซื้อ แต่ Use Case 5A ให้เจ้าของกดอนุมัติก่อนทุกครั้ง
\evidence{สไลด์ 9 และ Use Case Description 5A}
\end{q}

\begin{q}{อะไรที่ยังทำด้วยมือ}
\ans
วางแผนสินค้า จัดหา เจรจาโรงงาน ตัดสินใจเรื่องคุณภาพ อนุมัติการสั่งซื้อ กดพิมพ์ใบปะ และส่งมอบพัสดุให้ขนส่ง ส่วนลูกค้ายังสั่งและชำระเงินผ่าน Rakuten ตามเดิม
\end{q}

% ================================================================ คนที่ 2
\speaker{sB}{\speakerB}{Business Process, Use Case, diagram, ER}

\begin{q}{Use Case ที่เพิ่มจากการวิเคราะห์ CRUD มีอะไรบ้าง \must}
\ans
ตรวจคำสั่งซื้อ (2A) จับคู่กฎ SKU (2S) และตั้งกฎ SKU (3A) พิมพ์ใบปะ (9A) แยกจากการจัดรูปแบบ (6S) ส่งเลขติดตาม (5S) ตัดสต๊อก (7S) และลบข้อมูลเก่า (10A) ซึ่งเป็น Use Case เดียวที่มีการ Delete
\evidence{หัวข้อ 2.10 และ 2.12}
\end{q}

\begin{q}{ทำไม 1A และ 10A ไม่อยู่ใน Business Process}
\ans
1A เข้าสู่ระบบเป็นเงื่อนไขก่อนเริ่มทุกขั้นตอน ส่วน 10A เป็นงานดูแลระบบที่ทำเป็นระยะ ไม่ได้ผูกกับการรับส่งคำสั่งซื้อ จึงเขียนเหตุผลไว้ใต้ตารางจับคู่ BP กับ UC
\evidence{หัวข้อ 2.7}
\end{q}

\begin{q}{object ที่ติดกับ actor ใน Sequence Diagram คืออะไร \must}
\ans
เป็นหน้าจอทุกตัว เช่น LoginScreen หรือ ReorderDecisionScreen ไม่ใช้ชื่อ View และทุก diagram ต่อฐานข้อมูลผ่าน DBConnector เรามีสคริปต์ check\_diagrams.py ตรวจทุก diagram ว่า Admin ไม่ส่งข้อความตรงไปที่ controller และ object ใน Collaboration ตรงกับ Sequence
\end{q}

\begin{q}{เลข query อ่านอย่างไร}
\ans
Q ตามด้วยเลข Use Case และลำดับในขั้นตอน เช่น Q3A.1 คือ query แรกของ Use Case 3A เลขเดียวกันใช้ใน Use Case Description, Sequence Diagram, ER Diagram พร้อม Query, Query Table และคอมเมนต์ในโค้ด
\end{q}

\begin{q}{Class Diagram มีความสัมพันธ์แบบไหนบ้าง}
\ans
มีครบ 4 แบบ สืบทอดจาก BaseScreen, BaseController และ BaseRepository, association จากหน้าจอไปที่ Controller, dependency จาก Controller ไปที่ Repository และ aggregation ระหว่าง Repository กับ DBConnector
\evidence{รูป Class Diagram 7 รูปในหัวข้อ 2.14}
\end{q}

\begin{q}{State Diagram บอก SQL อย่างไร}
\ans
ทุกเส้นเขียนเป็น [เงื่อนไข] event / Use Case และกำกับเลข query ที่ UPDATE ค่า order\_status เช่น จากรอตรวจสอบคำสั่งซื้อไปรอจับคู่กฎ SKU ใช้ Q2A.4
\end{q}

\begin{q}{ER มี M:N ไหม total participation ดูตรงไหน}
\ans
ไม่มี M:N ทุก relationship เป็น 1:1 หรือ 1:N เส้นคู่คือ total participation เช่น ใบปะทุกใบต้องมี Order ส่วนเส้นเดี่ยวคือ partial เช่น Order ไม่จำเป็นต้องมีคำสั่งซื้อเติมสต๊อก
\evidence{หัวข้อ 2.18}
\end{q}

\begin{q}{ทำไม ER แบ่งเป็นหลายรูป}
\ans
ตาราง orders มี 34 คอลัมน์ ถ้ารวมทุก attribute ไว้ในรูปเดียวจะอ่านไม่ออกในกระดาษ A4 จึงมีรูปภาพรวมที่แสดงทุก entity และ relationship แล้วแยกรายละเอียด attribute เป็น 3 ส่วน
\end{q}

% ================================================================ คนที่ 3
\speaker{sC}{\speakerC}{SQL ฐานข้อมูล เว็บ การทดสอบ}

\begin{q}{SQL รันจริงหรือเปล่า ผลดูตรงไหน \must}
\ans
รันจริง เราสร้างฐานข้อมูล PostgreSQL 15 ตารางพร้อมข้อมูลตัวอย่าง แล้วรันทั้ง \num{63 คำสั่ง} ผ่านครบ แต่ละคำสั่งรันแยกจากข้อมูลตั้งต้นชุดเดียวกัน คำสั่งที่แก้ข้อมูลแสดงข้อมูลหลังรันด้วย
\evidence{หัวข้อ 2.25 และไฟล์ sql/schema.sql, seed.sql}
\end{q}

\begin{q}{ทำไม SQL ไม่ใช้ JOIN ซับซ้อน}
\ans
เกณฑ์ของวิชาให้ใช้ SQL ที่อ่านเข้าใจง่าย เราจึงแตกเป็นคำสั่งสั้น ๆ หลายข้อ เช่น 10A แยก query ของต้นทุน ใบปะ และคำสั่งซื้อเติมสต๊อกออกจากกัน
\end{q}

\begin{q}{ตอนนี้เว็บต่อฐานข้อมูลจริงหรือยัง}
\ans
หน้าจอทั้ง 12 หน้าใช้ข้อมูลตัวอย่างชุดเดียวกับฐานข้อมูลที่รัน SQL และเก็บใน localStorage ระหว่างพัฒนา โครงสร้างข้อมูลตรงกับ schema ทุกคอลัมน์ ขั้นต่อไปคือเปลี่ยนไปต่อ PostgreSQL จริงที่ไฟล์เดียว
\end{q}

\begin{q}{ทดสอบกรณีระบบภายนอกล่มอย่างไร \must}
\ans
มีหน้า /dev ที่ไม่อยู่ในเมนู ใช้จำลองข้อผิดพลาดตามทางเลือกของแต่ละ Use Case เช่น Rakuten เชื่อมต่อไม่ได้ หรือเครื่องพิมพ์ไม่พร้อม แล้วดูว่าระบบแสดงข้อความตรงกับ Use Case Description คำต่อคำ ทดสอบ \num{34 กรณี} ผ่านทุกกรณี
\evidence{หัวข้อ 2.26}
\end{q}

\begin{q}{ต้นทุนต่อหน่วยคำนวณอย่างไร}
\ans
Unit Cost = ((purchase\_price × exchange\_rate) + intl\_freight + duty\_fee) / order\_qty + marketplace\_fee + domestic\_shipping + rsl\_charge ตามสูตร Q4S.2 เช่น FSH-MLR-240 ได้ \num{418.61} ผลบันทึกใน cost\_calculation\_log ด้วย Q4S.3
\end{q}

\begin{q}{ถ้าจะขยายไป Amazon ต้องแก้อะไร}
\ans
เพิ่มค่าช่องทางขายและกฎ SKU ของช่องทางใหม่ในตาราง sku\_rules และ shipping\_rules ส่วน Use Case เดิมใช้ได้เพราะออกแบบให้ทุกตารางมี sales\_channel อยู่แล้ว
\end{q}

\end{document}
